\documentclass[11pt]{article}
\usepackage[T1]{fontenc}
\usepackage{lmodern,amsmath,amssymb,amsthm}
\usepackage[margin=1in]{geometry}
\usepackage[hidelinks]{hyperref}
\newtheorem{theorem}{Theorem}
\newtheorem{lemma}[theorem]{Lemma}
\newtheorem{corollary}[theorem]{Corollary}
\newcommand{\C}{\mathbb C}
\newcommand{\ord}{\operatorname{ord}_t}
\newcommand{\Nash}{\overline N}
\newcommand{\Cont}{\overline{\operatorname{Cont}}}
\newcommand{\A}{\mathcal A}
\renewcommand{\labelitemi}{\(\bullet\)}
\title{A free-chain counterexample to a Nash-intersection conjecture}
\author{Alec Kriebel\\
\small Independent researcher\\
\small ORCID: \href{https://orcid.org/0009-0001-9320-500X}{0009-0001-9320-500X}}
\date{22 October 2026\\\small Version 1.0.0 --- unrefereed preprint}
\begin{document}
\maketitle
\begin{abstract}
We compute the intersection of the closed Nash sets of the complete family
of prime divisors obtained by \(n\) successive free point blowups over
\((\C^2,0)\), after fixing a prefix of \(k\) centers, where \(2\leq k<n\).
In suitable coordinates the intersection is
\(\{\ord x\geq\lceil n/k\rceil,\ \ord y\geq n\}\).
For five blowups, fixing two or three centers gives the thresholds
\((3,5)\) or \((2,5)\). Neither is \(q(1,j)\) for positive integers \(q,j\).
These examples disprove the specified-reference-multiple assertion in
Conjecture~1 of Fern\'andez de Bobadilla, Pe Pereira and Popescu-Pampu,
under each contact and indexing convention explicitly treated here.
The source's bare equivalence symbol is undefined and its indexing is
inconsistent; no unique author-intended correction is asserted.
The intersections are themselves maximal divisorial sets for other
toric divisors. Thus the result does not disprove an assertion allowing
an arbitrary divisor, or solve the generalized Nash problem.
\end{abstract}

\section{The question and the precise negative answer}

For a prime exceptional divisor \(E\) over \((\C^2,0)\), let \(\Nash_E\)
be the Zariski closure of the nonconstant arcs whose lifts to its minimal
point-blowup model meet \(E\). We work throughout with closed sets of arcs,
equivalently reduced supports, in the centered arc space
\[
 \A=\{(x(t),y(t)):x,y\in t\C[[t]]\},\qquad \ord 0=+\infty.
\]
All cylinders below are defined by finitely many coefficient equations.
For \(q\geq1\), write \(\Cont^q(E)\) for the closure of arcs whose lifts
meet \(E\) with exact intersection multiplicity \(q\). We will also check
the convention using multiplicity at least \(q\).

Conjecture~1 in \cite[p.~31]{original}, and in the workshop account
\cite[p.~2679]{workshop}, proposes that an intersection of Nash sets
of divisors of the same combinatorial type with a prescribed common
prefix equals \(\Cont^q(E_p)\), where \(p\) is a chosen free point and
\(E_p\) is the divisor created by blowing it up. The surrounding
definition in \cite[p.~30]{original} defines \(E'\equiv_{\geq i}E\) by
equal combinatorial type and at least \(i\) common blown-up centers.
The conjecture itself prints a bare \(\equiv_{i_0}\), without defining
that symbol there. Remark~4.14 indexes the first-exceptional case in a
way that does not agree with a mechanical zero-based indexing of all
created centers. These are real limits on interpretation.

To avoid building a conclusion on an invented correction, we index
centers \(c_1=0,c_2,\ldots,c_n\), and let \(D_j\) be created at \(c_j\).
Continue the zero reference chain if \(D_j\) with \(j>n\) is mentioned.
A prefix of length \(k\) includes \(c_k\). The examples below treat
both fixing or not fixing the selected free \(c_3\), both common-prefix
and conditional exact-contact families, and every reference \(D_j\).
They are counterexamples to all the expressly listed readings, not
an assertion about an unspecified replacement of the printed conjecture.
The repository catalogue's union symbol is a transcription discrepancy:
both inspected primary formulations use an intersection.

\begin{theorem}\label{thm:main}
Let \(D_1,\ldots,D_n\) be the free chain following \(y=0\), with
\(n\geq3\), and let \(2\leq k<n\). Let \(\mathcal F_{n,k}\) be the
complete family of final prime divisors with this all-free combinatorial
type and with the first \(k\) centers fixed. Then
\[
 I_{n,k}:=\bigcap_{E'\in\mathcal F_{n,k}}\Nash_{E'}
   =\{(x,y)\in\A:\ord x\geq\lceil n/k\rceil,\ \ord y\geq n\}.
 \tag{1}\label{eq:intersection}
\]
The same intersection results if one restricts to divisors with exactly
\(k\) common centers with the zero reference chain.
For every \(j,q\geq1\),
\[
 \Cont^q(D_j)=\{\ord x\geq q,\ \ord y\geq jq\}.                 \tag{2}
 \label{eq:reference}
\]
This remains true for the closure of the at-least-\(q\) contact locus.
Consequently, for \(n=5\), neither \(I_{5,2}\) nor \(I_{5,3}\)
equals \(\Cont^q(D_j)\) for any positive integers \(q,j\).
\end{theorem}

The decisive assertion is the last conclusion: the required reference
divisor cannot be repaired by merely shifting its index. The general
formula also supplies positive boundary checks. In particular
\(I_{n,k}=\Cont^q(D_k)\) holds exactly when \(k\mid n\), with \(q=n/k\).

\section{The entire free family}

The source explicitly allows free as well as satellite prime divisors;
the free/satellite sequence with proximity data encodes combinatorial
type \cite[pp.~5--7]{original}. A point is free if it lies on the smooth
locus of the exceptional divisor. No restriction to final divisors of
minimal embedded resolutions of singular branches is imposed. Smooth
curvettas and the extra free blowups after resolving them are included.

\begin{lemma}\label{lem:family}
Every member of \(\mathcal F_{n,k}\), and no additional divisor, is
described by
\[
 P(x)=a_kx^k+\cdots+a_{n-1}x^{n-1},\qquad
 y=P(x)+x^n w,\qquad (a_k,\ldots,a_{n-1})\in\C^{n-k}.       \tag{3}
 \label{eq:family}
\]
Thus this is the full family, not a selected subfamily.
\end{lemma}

\begin{proof}
After \(j\) blowups along \(y=0\), the finite chart is \(y=x^j u_j\).
For \(j\geq2\), the newest exceptional divisor has equation \(x=0\);
its sole intersection with older exceptional components is the point
\(u_j=\infty\). Blowing up a finite point \(u_j=a\) gives the two charts
\[
 u_j=a+xu_{j+1},\qquad
 x=vh,\quad u_j=a+h.                                      \tag{4}
\]
In the second chart the new component is \(h=0\), the old component is
\(v=0\), and on the overlap \(u_{j+1}=1/v\).
The one point omitted from the first chart is therefore precisely the
satellite intersection, not an omitted free choice.

Starting with the fixed prefix, repeated substitution of the first
formula in (4) gives exactly (3). Conversely every finite coefficient
tuple chooses only free points. Different tuples first differing in
one coefficient choose a different center before creation of the final
divisor, and hence a different final divisor.

For completeness, the minimal sequence for one final prime divisor
is its chain of predecessors. A center on an older component away from
the newest component would make the intervening blowup extraneous,
so cannot occur in that minimal sequence. A satellite choice changes
the proximity data, and therefore the specified all-free type. This
excludes other apparent completions. Its weighted graph is the chain
with \(n-1\) weights \(-2\) and final weight \(-1\).
The initial fixed tangent allows \(x\) to be the parameter of every
smooth curvetta under discussion.
\end{proof}

For a nonzero prefix jet one first makes the polynomial coordinate
change subtracting that jet from \(y\); this gives the same statement
in regular coordinates. The counterexample needs only the displayed
zero reference. In the conditional exact-contact reading, exactly
\(k\) common centers means \(a_k\ne0\) in (3).

\section{Closures, including satellite limits}

\begin{lemma}\label{lem:cylinder}
For the divisor \(E_P\) in (3), put \(z=y-P(x)\). Then
\[
 \Nash_{E_P}=S_P:=\{\ord x\geq1,\ \ord z\geq n\}.           \tag{5}
\]
This includes limiting arcs with \(x=0\), the constant arc, and
lifts meeting a satellite point.
\end{lemma}

\begin{proof}
The coordinate change \((x,y)\mapsto(x,z)\) is polynomial and invertible.
The equations of \(S_P\) are just the vanishing of the first \(n-1\)
coefficients of \(z\). Its coefficient ring is a polynomial ring in the
remaining variables, so \(S_P\) is a reduced irreducible cylinder.
The nonempty open subset \(\ord x=1\) is dense. Every arc there has
\(w=z/x^n\in\C[[t]]\), so lifts in (3), meeting the newest component
transversely away from its older intersection. Hence \(S_P\) is contained
in the closure of the transverse locus.

For the reverse inclusion, the two charts covering the final component
have maps
\[
 (x,w)\longmapsto(x,z=x^n w),\qquad
 (v,h)\longmapsto(x=vh,\ z=v^{n-1}h^n).                    \tag{6}
\]
An arc centered on the new component has respectively \(\ord x\geq1\),
or \(\ord h\geq1\); in either case \(\ord x\geq1,\ord z\geq n\).
This proves containment even at its satellite intersection and hence (5).

One can check density directly without omitting degenerate arcs.
For any \((x_0,y_0)\in S_P\), set
\[
 x_\epsilon=x_0+\epsilon t,\qquad
 y_\epsilon=P(x_\epsilon)+y_0-P(x_0).
\]
The coefficients are polynomial in \(\epsilon\); the arc at zero is the
given one, and away from at most one value of \(\epsilon\) it lies in the
transverse locus. This handles \(x_0=0\) as well. A separate specialization
for each \(P\) suffices for membership in each closure.
\end{proof}

\begin{proof}[Proof of (2)]
For \(D_j\) the charts are (6) with \(P=0\) and \(n=j\).
The contact order in the first chart is \(\ord x\); in the second it is
\(\ord h\). Exact contact \(q\), or contact at least \(q\), therefore
implies \(\ord x\geq q,\ord y\geq jq\).
Conversely the cylinder with those two inequalities is irreducible, and
its nonempty dense open subset \(\ord x=q\) lifts in the first chart
with exact contact \(q\). Thus the closures for both conventions are
that cylinder. This argument works also for \(j=1\), whose two charts
are the usual first-blowup charts.
\end{proof}

\section{Intersections and uniform separation}

\begin{proof}[Proof of (1) and the exact-contact variant]
The members \(P=0\) and \(P=x^k\) force
\(\ord y\geq n\) and \(\ord x^k\geq n\). Since orders of nonzero
formal series are integers, the latter is
\(\ord x\geq\lceil n/k\rceil\).
Conversely, these inequalities imply \(\ord P(x)\geq n\) for every
polynomial in (3), because all its exponents are at least \(k\).
This proves (1) for the entire family.

For the exact-contact family, \(a_k\ne0\). Its two members \(P=x^k\)
and \(P=2x^k\) force \(\ord x^k\geq n\) by subtraction, and then
\(\ord y\geq n\). The converse is unchanged. This calculation treats
a possible reading; it does not assert that the source defined its bare
relation that way.
\end{proof}

Order cylinders with finite independent thresholds \((a,b)\) and
\((a',b')\) are equal only if \(a=a'\) and \(b=b'\).
For example, give one coordinate exactly its lower threshold and the
other arbitrarily high order, and interchange their roles.
Taking \(n=5\) gives
\[
 I_{5,2}=\{\ord x\geq3,\ord y\geq5\},\qquad
 I_{5,3}=\{\ord x\geq2,\ord y\geq5\}.                     \tag{7}
\]
Equality with (2) would require \((q,jq)=(3,5)\) or \((2,5)\).
Neither \(3j=5\) nor \(2j=5\) has an integer solution. This proves
the remaining conclusion of Theorem~\ref{thm:main} for every \(q,j\),
without a bounded search.

Here are explicit convergent polynomial separating arcs, also fixing
the direction of failure. Put \(r=3\) or \(2\).
If \(q<r\), the arc
\((t^q,t^{\max(5,jq)})\) lies in (2) and outside (7).
If \(q>r\), \((t^r,t^5)\) lies in (7) and outside (2).
If \(q=r\), use the latter arc when \(jr>5\), and use
\((t^r,t^{jr})\) when \(jr<5\). The remaining case \(jr=5\)
does not occur.

\begin{center}
\small
\begin{tabular}{p{.48\textwidth}
                p{.18\textwidth}
                p{.22\textwidth}}
Explicit normalization & Intersection & Reference tested\\ \hline
Last-shared-divisor convention, prefix \(c_1,c_2\)
 & \(I_{5,2}\) & \(D_2\)\\
Zero-based \(p_0=c_1,p_2=c_3\), contact at least \(2\)
 & \(I_{5,2}\) & \(D_3\)\\
Fix the selected free \(c_3\) as well
 & \(I_{5,3}\) & \(D_3\)\\
One-based \(p_3=c_3\), contact at least \(3\)
 & \(I_{5,3}\) & \(D_3\)\\
Conditional exact contact \(2\) or \(3\)
 & \(I_{5,2}\) or \(I_{5,3}\) & Every \(D_j\)
\end{tabular}
\end{center}
In particular \(c_3\) is genuinely free, and fixing it still leaves the
two parameters for \(c_4,c_5\). A three-blowup example can become a
singleton after fixing its final center; the five-blowup example avoids
that boundary issue. The raw transverse loci are not substituted for
their closures.

\section{Scope, prior framework and reproducibility}

The result concerns the \emph{specified reference multiple}. The
cylinders in (7) are themselves maximal divisorial sets for the primitive
toric valuations with weights \((3,5)\) and \((2,5)\). This follows from
the established toric description \cite[Example~2.10]{ishii}; adjacent
unimodular charts are respectively
\[
 (x,y)=(u^3v,u^5v^2),\qquad (x,y)=(u^2v,u^5v^3).
\]
Their exponent matrices have determinant \(1\). The inverse monomials are
\[
 (u,v)=(x^2/y,y^3/x^5),\qquad (u,v)=(x^3/y,y^2/x^5).
\]
The dense exact-order loci have \(\ord u=1\) and \(v\) a unit.
Thus allowing a different arbitrary prime divisor removes this
particular obstruction. No negative solution of the generalized Nash
problem, no statement about Conjecture~2, and no new general theory of
maximal divisorial sets is claimed.

We also distinguish closed sets from nonreduced intersections. For
\(n=5,k=2\), the two-member coefficient ideal after eliminating the first
four coefficients of \(y\) is
\[
 (x_1^2,\ 2x_1x_2,\ x_2^2+2x_1x_3),
\]
whose radical is \((x_1,x_2)\), although \(x_1\) is not in that ideal.
For \(k=3\) the corresponding ideal is
\((x_1^3,3x_1^2x_2)\), with radical \((x_1)\).
Equation (7) states the reduced supports, not equality of these
nonreduced schemes. Divisible cases \(n=6,k=2,3\) satisfy the reference
equality, and fixing the final center gives the singleton \(D_n\).
Varying all tangents for an all-free length-\(n\) type gives
\(\{\ord x,\ord y\geq n\}=\Cont^n(D_1)\), consistent with the
first-exceptional specialization rather than refuting it.

The standard maximal-divisorial/contact framework is due to prior work,
including \cite{ishii}; combinatorial invariance and finite cylindrical
intersections are proved in \cite{original}. The contribution here is
the explicit complete free-family calculation and its reference-ray
obstruction. The source's authors still discuss their two conjectures
as conjectures in \cite[Section~8]{survey}. An extensive but bounded
primary-source investigation is documented in the accompanying
\texttt{SOURCE\_AUDIT.md}. It located no authenticated earlier
publication of this exact negative calculation. This is not a certificate
of worldwide firstness, continuous openness, or completeness of the
literature search. The exact claim is tied to the inspected
arXiv1501.06898v3 and workshop2016 bodies. The publisher-typeset full text
was not accessed; the original paper's DOI is a reference, not this
counterexample's publication DOI. No author-approved erratum, intended
new hypothesis, or unique resolution of its notation is asserted.

The source package includes an independently written standard-library
checker, expected output, editable TeX, the reviewed PDF, source/provenance
notes and reproduction instructions. Exact finite computations challenge
coefficient substitutions, two-chart formulas, degeneracies, all listed
finite witness cases, and boundary mutations. They supplement the
universal proofs above; they do not certify arbitrary complex parameters,
priority, or infinite-arc geometry. The package builder reproduces
archives from the frozen PDF; it does not pretend to typeset that PDF.

\paragraph{AI and review disclosure.}
AI tools were extensively involved in mathematical research, proof
development and reconstruction, symbolic and exact computation,
literature/source investigation, manuscript preparation, editing and
adversarial auditing. Alec Kriebel is the named human author.
This manuscript is an unrefereed preprint and has not undergone formal
independent human peer review or journal acceptance. Internal AI audits
and exact supplementary computations do not provide those certifications
or proof-assistant verification.

\begin{thebibliography}{9}\small
\bibitem{original}
J. Fern\'andez de Bobadilla, M. Pe Pereira and P. Popescu-Pampu,
\emph{On the generalized Nash problem for smooth germs and adjacencies
of curve singularities}, Adv. Math. \textbf{320} (2017), 1269--1317,
\href{https://doi.org/10.1016/j.aim.2017.09.035}{doi:10.1016/j.aim.2017.09.035}.
Inspected author version: \href{https://arxiv.org/abs/1501.06898v3}
{arXiv:1501.06898v3} (5 October 2017).
\bibitem{workshop}
M. Pe Pereira (joint work with J. Fern\'andez de Bobadilla and
P. Popescu-Pampu), \emph{The arcspace of \(\C^2\) and adjacencies of
plane curves}, in \emph{Singularities}, Oberwolfach Reports
\textbf{46/2016}, contribution pp.~2678--2681, Conjecture~1 on p.~2679.
Primary report:
\url{https://ems.press/content/serial-article-files/46651}.
\bibitem{ishii}
S. Ishii, \emph{Maximal divisorial sets in arc spaces}, in
\emph{Algebraic Geometry in East Asia---Hanoi 2005}, Adv. Stud. Pure
Math. \textbf{50} (2008), 237--249,
\href{https://doi.org/10.2969/aspm/05010237}{doi:10.2969/aspm/05010237}.
Inspected author version: \href{https://arxiv.org/abs/math/0612185}
{arXiv:math/0612185}.
\bibitem{survey}
J. Fern\'andez de Bobadilla and M. Pe Pereira,
\emph{The Nash problem from a geometric and topological perspective},
\href{https://arxiv.org/abs/1805.01418}{arXiv:1805.01418} (2018),
Section~8, pp.~16--17; cited body is this arXiv edition.
\end{thebibliography}
\end{document}
